\documentclass[10pt,a4paper]{amsart}
\usepackage{amsmath,amssymb,amsthm,graphicx}
\usepackage{verbatim}
\newenvironment{reference}{\comment}{\endcomment}
%\newenvironment{reference}{}{}
\newenvironment{slogan}{\comment}{\endcomment}
\newenvironment{history}{\comment}{\endcomment}
% For commutative diagrams we use Xy-pic
\usepackage[all]{xy}

% We use 2cell for 2-commutative diagrams.
\xyoption{2cell}
% Theorem environments.
%
\theoremstyle{plain}
\newtheorem{theorem}[subsection]{Theorem}
\newtheorem{proposition}[subsection]{Proposition}
\newtheorem{lemma}[subsection]{Lemma}

\theoremstyle{definition}
\newtheorem{definition}[subsection]{Definition}
\newtheorem{example}[subsection]{Example}
\newtheorem{exercise}[subsection]{Exercise}
\newtheorem{situation}[subsection]{Situation}

\theoremstyle{remark}
\newtheorem{remark}[subsection]{Remark}
\newtheorem{remarks}[subsection]{Remarks}

\numberwithin{equation}{subsection}

% Macros
%
\def\lim{\mathop{\mathrm{lim}}\nolimits}
\def\colim{\mathop{\mathrm{colim}}\nolimits}
\def\Spec{\mathop{\mathrm{Spec}}}
\def\Hom{\mathop{\mathrm{Hom}}\nolimits}
\def\Ext{\mathop{\mathrm{Ext}}\nolimits}
\def\SheafHom{\mathop{\mathcal{H}\!\mathit{om}}\nolimits}
\def\SheafExt{\mathop{\mathcal{E}\!\mathit{xt}}\nolimits}
\def\Sch{\mathit{Sch}}
\def\Mor{\mathop{\mathrm{Mor}}\nolimits}
\def\Ob{\mathop{\mathrm{Ob}}\nolimits}
\def\Sh{\mathop{\mathit{Sh}}\nolimits}
\def\NL{\mathop{N\!L}\nolimits}
\def\CH{\mathop{\mathrm{CH}}\nolimits}
\def\proetale{{pro\text{-}\acute{e}tale}}
\def\etale{{\acute{e}tale}}
\def\QCoh{\mathit{QCoh}}
\def\Ker{\mathop{\mathrm{Ker}}}
\def\Im{\mathop{\mathrm{Im}}}
\def\Coker{\mathop{\mathrm{Coker}}}
\def\Coim{\mathop{\mathrm{Coim}}}

% Boxtimes
%
\DeclareMathSymbol{\boxtimes}{\mathbin}{AMSa}{"02}

%
% Macros for moduli stacks/spaces
%
\def\QCohstack{\mathcal{QC}\!\mathit{oh}}
\def\Cohstack{\mathcal{C}\!\mathit{oh}}
\def\Spacesstack{\mathcal{S}\!\mathit{paces}}
\def\Quotfunctor{\mathrm{Quot}}
\def\Hilbfunctor{\mathrm{Hilb}}
\def\Curvesstack{\mathcal{C}\!\mathit{urves}}
\def\Polarizedstack{\mathcal{P}\!\mathit{olarized}}
\def\Complexesstack{\mathcal{C}\!\mathit{omplexes}}
% \Pic is the operator that assigns to X its picard group, usage \Pic(X)
% \Picardstack_{X/B} denotes the Picard stack of X over B
% \Picardfunctor_{X/B} denotes the Picard functor of X over B
\def\Pic{\mathop{\mathrm{Pic}}\nolimits}
\def\Picardstack{\mathcal{P}\!\mathit{ic}}
\def\Picardfunctor{\mathrm{Pic}}
\def\Deformationcategory{\mathcal{D}\!\mathit{ef}}
\usepackage{cleveref}
\crefname{originallemma}{original lemma}{original lemmas}
\Crefname{originallemma}{Original lemma}{Original lemmas}
\makeatletter
\newlabel{coherent-lemma-quasi-compact-h1-zero-covering}{{01XF}{source}{}{sourcepointer-1}{}}
\newlabel{coherent-lemma-quasi-compact-h1-zero-covering@cref}{{[originallemma][0][]01XF}{[1][0][]source}{}{}{}}
\makeatother
\crefname{originallemma}{original lemma}{original lemmas}
\Crefname{originallemma}{Original lemma}{Original lemmas}
\makeatletter
\newlabel{lemma-flat-base-change-cohomology}{{073K}{source}{}{sourcepointer-2}{}}
\newlabel{lemma-flat-base-change-cohomology@cref}{{[originallemma][0][]073K}{[1][0][]source}{}{}{}}
\makeatother
\crefname{originallemma}{original lemma}{original lemmas}
\Crefname{originallemma}{Original lemma}{Original lemmas}
\makeatletter
\newlabel{lemma-vanishing-injective}{{07V4}{source}{}{sourcepointer-3}{}}
\newlabel{lemma-vanishing-injective@cref}{{[originallemma][0][]07V4}{[1][0][]source}{}{}{}}
\makeatother
\crefname{originallemma}{original lemma}{original lemmas}
\Crefname{originallemma}{Original lemma}{Original lemmas}
\makeatletter
\newlabel{lemma-vanishing-separated}{{07V5}{source}{}{sourcepointer-4}{}}
\newlabel{lemma-vanishing-separated@cref}{{[originallemma][0][]07V5}{[1][0][]source}{}{}{}}
\makeatother
\crefname{originalreference}{original reference}{original references}
\Crefname{originalreference}{Original reference}{Original references}
\makeatletter
\newlabel{situation-vanishing}{{07UY}{source}{}{sourcepointer-5}{}}
\newlabel{situation-vanishing@cref}{{[originalreference][0][]07UY}{[1][0][]source}{}{}{}}
\makeatother
\crefname{originallemma}{original lemma}{original lemmas}
\Crefname{originallemma}{Original lemma}{Original lemmas}
\makeatletter
\newlabel{spaces-morphisms-lemma-locally-quasi-finite}{{06RW}{source}{}{sourcepointer-6}{}}
\newlabel{spaces-morphisms-lemma-locally-quasi-finite@cref}{{[originallemma][0][]06RW}{[1][0][]source}{}{}{}}
\makeatother
\crefname{originallemma}{original lemma}{original lemmas}
\Crefname{originallemma}{Original lemma}{Original lemmas}
\makeatletter
\newlabel{spaces-morphisms-lemma-locally-quasi-finite-separated-representable}{{0418}{source}{}{sourcepointer-7}{}}
\newlabel{spaces-morphisms-lemma-locally-quasi-finite-separated-representable@cref}{{[originallemma][0][]0418}{[1][0][]source}{}{}{}}
\makeatother
\title{Serre's criterion for affineness in the setting of algebraic spaces, by a presentation of the structure sheaf of an affine etale cover}
\author{The Stacks Project Authors}
\date{stacks-project:spaces-cohomology.tex:local-snapshot-20261012; GNU FDL 1.2 or later}
\begin{document}
\maketitle
\noindent\textit{Editorial note.} All delivered mathematical text is the author's own, copied from the pinned Stacks Project chapter file spaces-cohomology.tex; no statement or proof step was rewritten, reordered or removed.\par\medskip
\noindent\textit{Editorial note.} Preamble provenance: the pinned chapter file opens with \textbackslash{}input\{preamble\} and carries no preamble text of its own, so the project's single Stacks-wide preamble staged at sources/stacks-spaces-cohomology/preamble.tex is excerpted verbatim as three spans, lines 8-12 (verbatim package and the project's own reference/slogan/history display environments), lines 14-18 (Xy-pic setup) and lines 154-217 (theorem declarations and mathematical macros). The wrapper class and the amsmath/amssymb/amsthm/graphicx base packages are supplied by the assembler. No macro was invented or redefined.\par\medskip
\noindent\textit{Editorial note.} Same-chapter pointers are located at their \textbackslash{}label lines inside the pinned chapter file; cross-chapter pointers are located in the project's own tag list staged in the snapshot (sources/stacks-chow/tags.txt), at the verbatim lines carrying their canonical tags.\par\medskip
\noindent\textit{Allowed standard prerequisites (editorial).}\begin{itemize}
\item The vanishing situation is separated: the fibre product of an affine etale cover is a closed subscheme of the product of the base with itself. Same chapter spaces-cohomology.tex, label lemma-vanishing-separated (invoked at proof line 3173). It makes \$R = U \textbackslash{}times\_X U\$ a closed subscheme of \$\textbackslash{}Spec(B \textbackslash{}otimes\_A B)\$, hence affine, which is what allows the presentation of the algebra \$C\$.
\item The base change identification of the pullback of the pushforward of the structure sheaf. Same chapter spaces-cohomology.tex, label lemma-flat-base-change-cohomology (invoked at proof line 3191). It identifies \$\textbackslash{}varphi\textasciicircum{}*\textbackslash{}varphi\_*\textbackslash{}mathcal\{O\}\_U\$ with \$t\_*\textbackslash{}mathcal\{O\}\_R\$, so that surjectivity of the map on the cover becomes the generation condition \$s(g\_i)\$ over \$t\$.
\item The injectivity consequence of the vanishing situation, and representability of a separated locally quasi-finite morphism. Same chapter spaces-cohomology.tex, label lemma-vanishing-injective (invoked at proof line 3201) and Morphisms of Spaces, labels spaces-morphisms-lemma-locally-quasi-finite and spaces-morphisms-lemma-locally-quasi-finite-separated-representable (invoked at proof lines 3203 and 3206). Together with separatedness and finite type they make \$X \textbackslash{}to \textbackslash{}Spec(A)\$ representable, so that \$X\$ is a scheme.
\item The affine criterion for a quasi-compact separated scheme on which \$H\textasciicircum{}1\$ of every quasi-coherent module vanishes. Cohomology of Schemes, label coherent-lemma-quasi-compact-h1-zero-covering (invoked at proof line 3208). It finishes the argument by showing that the scheme \$X\$ is affine, hence equal to the base.
\end{itemize}\medskip
\ifdefined\hypertarget\hypertarget{sourcepointer-1}{}\fi
\noindent\textit{Source reference.} Original lemma 01XF: sources/stacks-chow/tags.txt, source lines 2516--2516. Navigation pointer for \textbackslash{}ref\{coherent-lemma-quasi-compact-h1-zero-covering\} used inside the counted span (line 3209); the argument that consumes it is retained in the span.\par\medskip
\ifdefined\hypertarget\hypertarget{sourcepointer-2}{}\fi
\noindent\textit{Source reference.} Original lemma 073K: sources/stacks-spaces-cohomology/pinned-original-source.tex, source lines 1922--1922. Navigation pointer for \textbackslash{}ref\{lemma-flat-base-change-cohomology\} used inside the counted span (line 3191); the argument that consumes it is retained in the span.\par\medskip
\ifdefined\hypertarget\hypertarget{sourcepointer-3}{}\fi
\noindent\textit{Source reference.} Original lemma 07V4: sources/stacks-spaces-cohomology/pinned-original-source.tex, source lines 3097--3097. Navigation pointer for \textbackslash{}ref\{lemma-vanishing-injective\} used inside the counted span (line 3201); the argument that consumes it is retained in the span.\par\medskip
\ifdefined\hypertarget\hypertarget{sourcepointer-4}{}\fi
\noindent\textit{Source reference.} Original lemma 07V5: sources/stacks-spaces-cohomology/pinned-original-source.tex, source lines 3132--3132. Navigation pointer for \textbackslash{}ref\{lemma-vanishing-separated\} used inside the counted span (line 3173); the argument that consumes it is retained in the span.\par\medskip
\ifdefined\hypertarget\hypertarget{sourcepointer-5}{}\fi
\noindent\textit{Source reference.} Original reference 07UY: sources/stacks-spaces-cohomology/pinned-original-source.tex, source lines 2991--2991. Navigation pointer for \textbackslash{}ref\{situation-vanishing\} used inside the counted span (line 3167); the argument that consumes it is retained in the span.\par\medskip
\ifdefined\hypertarget\hypertarget{sourcepointer-6}{}\fi
\noindent\textit{Source reference.} Original lemma 06RW: sources/stacks-chow/tags.txt, source lines 8655--8655. Navigation pointer for \textbackslash{}ref\{spaces-morphisms-lemma-locally-quasi-finite\} used inside the counted span (line 3203); the argument that consumes it is retained in the span.\par\medskip
\ifdefined\hypertarget\hypertarget{sourcepointer-7}{}\fi
\noindent\textit{Source reference.} Original lemma 0418: sources/stacks-chow/tags.txt, source lines 5283--5283. Navigation pointer for \textbackslash{}ref\{spaces-morphisms-lemma-locally-quasi-finite-separated-representable\} used inside the counted span (line 3206); the argument that consumes it is retained in the span.\par\medskip
\begin{proposition}
\label{proposition-vanishing-affine}
\begin{slogan}
Serre's criterion for affineness in the setting of algebraic spaces.
\end{slogan}
A quasi-compact and quasi-separated algebraic space is affine
if and only if all higher cohomology groups of quasi-coherent sheaves
vanish. More precisely, any algebraic space as in
Situation \ref{situation-vanishing} is an affine scheme.
\end{proposition}

\begin{proof}
Choose an affine scheme $U = \Spec(B)$ and a surjective \'etale
morphism $\varphi : U \to X$. Set $R = U \times_X U$. As $p$ is separated
(Lemma \ref{lemma-vanishing-separated}) we see that $R$ is a
closed subscheme of $U \times_{\Spec(A)} U = \Spec(B \otimes_A B)$.
Hence $R = \Spec(C)$ is affine too and the ring map
$$
B \otimes_A B \longrightarrow C
$$
is surjective. Let us denote the two maps $s, t : B \to C$ as usual. Pick
$g_1, \ldots, g_m \in B$ such that $s(g_1), \ldots, s(g_m)$ generate $C$
over $t : B \to C$ (which is possible as $t : B \to C$ is of finite
presentation and the displayed map is surjective). Then $g_1, \ldots, g_m$
give global sections of $\varphi_*\mathcal{O}_U$ and the map
$$
\mathcal{O}_X[z_1, \ldots, z_n] \longrightarrow \varphi_*\mathcal{O}_U,
\quad
z_j \longmapsto g_j
$$
is surjective: you can check this by restricting to $U$.
Namely, $\varphi^*\varphi_*\mathcal{O}_U = t_*\mathcal{O}_R$
(by Lemma \ref{lemma-flat-base-change-cohomology})
hence you get exactly the condition that $s(g_i)$ generate $C$
over $t : B \to C$. By the vanishing of $H^1$ of the kernel we see that
$$
\Gamma(X, \mathcal{O}_X[x_1, \ldots, x_n]) =
A[x_1, \ldots, x_n] \longrightarrow
\Gamma(X, \varphi_*\mathcal{O}_U) = \Gamma(U, \mathcal{O}_U) = B
$$
is surjective. Thus we conclude that $B$ is a finite type $A$-algebra.
Hence $X \to \Spec(A)$ is of finite type and separated.
By Lemma \ref{lemma-vanishing-injective}
and
Morphisms of Spaces, Lemma \ref{spaces-morphisms-lemma-locally-quasi-finite}
it is also locally quasi-finite. Hence $X \to \Spec(A)$ is representable by
Morphisms of Spaces, Lemma
\ref{spaces-morphisms-lemma-locally-quasi-finite-separated-representable}
and $X$ is a scheme. Finally $X$ is affine, hence equal to $\Spec(A)$,
by an application of Cohomology of Schemes, Lemma
\ref{coherent-lemma-quasi-compact-h1-zero-covering}.
\end{proof}

\end{document}
